\documentclass[11pt]{amsart}

\usepackage[margin=1in]{geometry}
\usepackage{amsmath,amssymb,mathtools}
\usepackage{array,booktabs}
\usepackage[T1]{fontenc}
\usepackage{lmodern}
\usepackage{microtype}
\usepackage{xcolor}
\usepackage{xurl}
\usepackage[colorlinks=true,
            linkcolor=blue!55!black,
            citecolor=blue!55!black,
            urlcolor=blue!55!black]{hyperref}

\newtheorem{theorem}{Theorem}[section]
\newtheorem{proposition}[theorem]{Proposition}
\newtheorem{lemma}[theorem]{Lemma}
\newtheorem{corollary}[theorem]{Corollary}
\theoremstyle{definition}
\newtheorem{definition}[theorem]{Definition}
\theoremstyle{remark}
\newtheorem{remark}[theorem]{Remark}

\newcommand{\Z}{\mathbf Z}
\newcommand{\F}{\mathbf F}
\newcommand{\Ga}{\mathbb G_a}
\newcommand{\Gm}{\mathbb G_m}
\newcommand{\Spec}{\operatorname{Spec}}
\newcommand{\Soc}{\operatorname{Soc}}
\newcommand{\Ann}{\operatorname{Ann}}
\newcommand{\Der}{\operatorname{Der}}
\newcommand{\Hom}{\operatorname{Hom}}
\newcommand{\gr}{\operatorname{gr}}
\newcommand{\length}{\operatorname{length}}
\newcommand{\eps}{\varepsilon}
\newcommand{\m}{\mathfrak m}
\newcommand{\n}{\mathfrak n}
\newcommand{\ot}{\otimes}
\newcommand{\id}{\mathrm{id}}
\newcommand{\unit}{\mathrm e}

\title[A fourth-power carry in rank four]
      {A Fourth-Power Carry in a Finite Flat\\
       Group Scheme of Rank Four}
\author{}
\date{}
\subjclass[2020]{14L15, 16T05}
\keywords{finite flat group scheme, Hopf algebra, power word,
Artin complete intersection}

\begin{document}

\begin{abstract}
The analogue of Lagrange's theorem for a finite locally free group scheme
asks whether its rank annihilates the power word.  We give an explicit
negative example in rank four.  Its base is the length-nine local Artin
complete intersection
\[
 C=\Z[r,s]/(r^3+2,\ s^3+2,\ rs^2),
\]
and its coordinate algebra is free on \(1,U,V,UV\).  Its special fiber is
\(\alpha _2\times\alpha _2\), and every geometric fiber is obtained from
this one by scalar extension.  The total group scheme is nevertheless
noncommutative and its fourth-power word is nontrivial; its eighth-power
word is trivial.

The construction is governed by a simple mechanism.  A group-like element
of the form \(\lambda=1+\theta\), with \(\theta^2=0\), leaves the geometric
sum carry
\[
 1+\lambda+\lambda^2+\lambda^3=2\theta
\]
in characteristic four.  Two elementary rank-two cancellation patterns,
joined by one mixed bridge, realize this mechanism inside a closed
submonoid of \(\Ga^2\rtimes\Gm\).  The carry becomes the socle class
\(2sUV\); periodicity then supplies the inverse and turns the submonoid into
a group scheme.  It may also be viewed as a square-zero deformation of an
exponent-four group: the failure to lift the equation \([4]=\unit\) is the
nonzero cotangent-theoretic derivation \(dU\mapsto2sUV\), \(dV\mapsto0\).
Thus the failure is entirely relative: it lives in the last infinitesimal
layer and is invisible on every geometric fiber.
\end{abstract}

\maketitle

\section{Introduction}

If a finite group has order \(n\), then every element has \(n\)-th power
equal to the identity.  The corresponding question for a finite locally
free group scheme \(G\to S\) of rank \(n\) is whether the power word
\[
 [n]\colon G\longrightarrow G,\qquad g\longmapsto g^n,
\]
is the constant unit word.  This question is recorded in SGA~3; see
\cite[Exp.~VIII, Rem.~7.3.1]{SGA3II} and the account in
\cite{SchoofSurvey}.  When \(G\) is noncommutative, the power word need not
be a homomorphism, but it remains a perfectly well-defined morphism of
schemes.  On a coordinate Hopf algebra \(A\), killedness by \(n\) means
\[
 [n]^\#=\eta\eps\colon A\longrightarrow A.
\]

The answer over a field is proved in
\cite[Exp.~VIIA, Prop.~8.5]{SGA3I}; it follows that the answer is also
positive over a reduced base.  Deligne proved it for commutative finite
locally free group schemes over an arbitrary base.  His argument first
appeared in \cite[\S1]{OortTate}; later expositions include
\cite[\S3.8]{TateFiniteFlat}, \cite[\S3.3]{Stix}, and
\cite[\S3]{SchoofSurvey}.  Any counterexample must therefore use both
nilpotents in the base and a noncommutative group law.

There are also strong positive results in the remaining, genuinely
relative setting.  Let \(R\) be a local Artin ring with residue
characteristic \(p>0\) and maximal ideal \(\mathfrak m_R\).  Schoof proved
that every finite flat group scheme over \(R\) is killed by its order if
\begin{equation}\label{eq:Schoof-range}
 \mathfrak m_R^p=p\mathfrak m_R=0;
\end{equation}
in particular this applies when \(\mathfrak m_R^2=0\)
\cite{Schoof2001,SchoofSurvey}.  Schoof also proved, without a restriction
on the length of \(R\), killedness for deformations of a distinguished
noncommutative matrix group in characteristic \(p\).  Torti extended this
deformation-theoretic result to an explicit family
\(G_\lambda\simeq\alpha_p\rtimes_\lambda\mu_{p^m}\) of noncommutative
special fibers \cite{Torti}.  These results, together
with the original references, are surveyed in
\cite{SchoofSurvey,Torti}.

The example below lies just beyond these positive ranges.  Its residue
characteristic is \(2\), but its maximal ideal satisfies
\[
 \mathfrak m^2\ne0,
 \qquad
 2\mathfrak m=(2s)\ne0,
\]
so \eqref{eq:Schoof-range} fails in both ways.  Moreover, its special fiber
is the commutative group \(\alpha_2\times\alpha_2\), rather than one of the
noncommutative semidirect products treated in the deformation results
above.  Noncommutativity itself is created by the nilpotent thickening.
Thus both features excluded by the classical theorems will be visible,
rather than incidental, in the construction.

There is a basic reason that the relative question can behave differently
from its fiberwise analogue.  A morphism over a nonreduced base may vanish
on every geometric fiber and still survive in a nilpotent direction.  The
example below makes this possibility concrete in a compact form: its
fourth-power defect is a generator of the socle of the total
coordinate algebra.

Here is the complete statement.

\begin{theorem}\label{thm:main}
Let
\begin{equation}\label{eq:C-intro}
 C=\Z[r,s]/(r^3+2,\ s^3+2,\ rs^2),
\end{equation}
and define
\begin{equation}\label{eq:A-intro}
 A=C[U,V]/(U^2-r^2U+rsV,\ V^2-s^2V).
\end{equation}
Put \(W=UV\) and
\begin{equation}\label{eq:lambda-intro}
 \lambda=(1+rU)(1+sV).
\end{equation}
Then the following assertions hold.
\begin{enumerate}
\item The ring \(C\) is a local Artin complete intersection of length nine,
with residue field \(\F_2\), characteristic four, and
\[
 \Soc(C)=\F_2(2s).
\]
\item The algebra \(A\) is free of rank four over \(C\), with basis
\[
 1,\qquad U,\qquad V,\qquad W.
\]
It is a Hopf algebra with counit \(\eps(U)=\eps(V)=0\) and coproduct
\begin{equation}\label{eq:Delta-intro}
 \Delta(U)=U\ot1+\lambda\ot U,
 \qquad
 \Delta(V)=V\ot\lambda+1\ot V.
\end{equation}
\item For \(G=\Spec A\), the special fiber is
\(\alpha _2\times\alpha _2\); consequently every geometric fiber is its
scalar extension to an algebraic closure.  The group scheme \(G\) itself
is noncommutative.
\item The power words satisfy
\begin{equation}\label{eq:powers-intro}
 [4]^\#(U)=2sW\ne0,
 \qquad
 [4]^\#(V)=[4]^\#(W)=0,
 \qquad
 [8]=\unit.
\end{equation}
In fact \(2sW\) spans the socle of the local Artin ring \(A\).
\end{enumerate}
In particular, \(G\) is a finite locally free group scheme of rank four
which is not killed by four.
\end{theorem}

The formulas in the theorem were not found by guessing a large Hopf table.
They can be reconstructed from one elementary observation.  If \(\lambda\)
is group-like and \(U\) is \(\lambda\)-skew primitive, then
\[
 [n]^\#(U)=(1+\lambda+\cdots+\lambda^{n-1})U.
\]
If \(\lambda=1+\theta\), \(\theta^2=0\), and \(4=0\), the four-term
geometric sum is \(2\theta\), not necessarily zero.  The problem is
therefore to build a rank-four algebra in which \(2\theta U\) survives.

The rest of the paper follows this reverse-engineering argument.  We first
isolate the geometric-sum carry.  We then explain how two rank-two Hopf
patterns and one bridge produce exactly the three equations in
\eqref{eq:C-intro}.  Next we study the geometry of the base and realize the
construction as a closed submonoid of an elementary matrix group.  The
eighth-power calculation supplies inverses automatically.  We then
identify the fourth-power carry as the obstruction to lifting a power
relation across a square-zero extension.  Finally, we describe the special
fiber, noncommutativity, a nonnormal rank-two subgroup, and the
precise---but deliberately limited---sense in which the three base
equations are economical.

\section{The design principle: skew primitives and geometric sums}

Let \(R\) be a commutative ring and let \(M=\Spec B\) be an affine monoid
scheme over \(R\).  We use the same terminology as for Hopf algebras: an
element \(l\in B\) is \emph{group-like} if
\(\Delta(l)=l\ot l\) and \(\eps(l)=1\), and \(x\in B\) is left
\(l\)-\emph{skew primitive} if
\[
 \Delta(x)=x\ot1+l\ot x.
\]
No inverse is needed for these notions.

\begin{remark}[The geometric dictionary]\label{rem:geometric-dictionary}
A function \(l\in B\) is the same thing as a morphism of schemes
\[
 l\colon M\longrightarrow\mathbb A^1_R.
\]
The two group-like identities say exactly that this morphism respects
multiplication and the unit, where the target is given its multiplicative
monoid structure.  Thus a group-like element is a character
\[
 M\longrightarrow\mathbb A^1_{\mathrm{mult}}.
\]
If \(M\) is a group scheme, then \(l\) is automatically invertible, with
inverse \(S(l)\), and the map lands in \(\Gm\).  In that case a group-like
element is literally a homomorphism of group schemes
\[
 l\colon M\longrightarrow\Gm.
\]

A skew-primitive function is almost, but not quite, a map to \(\Ga\).
When \(l\) is invertible, the pair \((x,l)\) is a homomorphism to the
affine group
\begin{equation}\label{eq:affine-group}
 \operatorname{Aff}_1=\Ga\rtimes\Gm
 =
 \left\{
 \begin{pmatrix}z&u\\0&1\end{pmatrix}
 \right\},
 \qquad
 (u,z)(u',z')=(u+zu',zz').
\end{equation}
Indeed, on points the two coproduct identities are precisely
\begin{equation}\label{eq:skew-primitive-on-points}
 l(gh)=l(g)l(h),
 \qquad
 x(gh)=x(g)+l(g)x(h),
\end{equation}
which say that
\[
 g\longmapsto
 \begin{pmatrix}l(g)&x(g)\\0&1\end{pmatrix}
\]
respects multiplication.  Equivalently, \(x\) is a crossed homomorphism,
or \(1\)-cocycle, for the one-dimensional action defined by the character
\(l\).  When \(l=1\), the twisting disappears:
\(\Delta(x)=x\ot1+1\ot x\), and \(x\colon M\to\Ga\) is an ordinary
homomorphism.

The right skew-primitive identity
\[
 \Delta(y)=y\ot l+1\ot y
\]
has the same meaning using lower-triangular matrices:
\[
 g\longmapsto
 \begin{pmatrix}l(g)&0\\y(g)&1\end{pmatrix},
 \qquad
 y(gh)=y(g)l(h)+y(h).
\]
Without the invertibility of \(l\), the same statements define a map to
the affine-transformation \emph{monoid}
\(\Ga\rtimes\mathbb A^1_{\mathrm{mult}}\).  In the construction below,
\(\lambda\) is a unit, so the target is already the affine group.
\end{remark}

\begin{lemma}[The norm formula]\label{lem:norm}
Suppose \(l\in B\) is group-like.  If \(x\) is left \(l\)-skew primitive,
then
\begin{equation}\label{eq:left-norm}
 [n]^\#(x)=N_n(l)x,
 \qquad
 N_n(T)=1+T+\cdots+T^{n-1}.
\end{equation}
If instead
\(\Delta(y)=y\ot l+1\ot y\), then
\begin{equation}\label{eq:right-norm}
 [n]^\#(y)=yN_n(l).
\end{equation}
\end{lemma}

\begin{proof}
In the affine-transformation monoid (and hence in the group
\eqref{eq:affine-group} when \(l\) is invertible), ordinary matrix
multiplication gives
\[
 (u,z)^n=\bigl(N_n(z)u,z^n\bigr).
\]
Pulling the first coordinate back along the homomorphism \((x,l)\) gives
\eqref{eq:left-norm}.  Equivalently, iterating the coproduct gives
\[
 x\ot1\ot\cdots\ot1
 +l\ot x\ot1\ot\cdots\ot1
 +\cdots+l\ot\cdots\ot l\ot x.
\]
Multiplying the tensor factors gives \eqref{eq:left-norm}.  The right
skew-primitive formula is the same calculation with the lower-triangular
matrices above.
\end{proof}

The identity which drives the example is now immediate.

\begin{lemma}[The characteristic-four carry]\label{lem:carry}
Assume \(4=0\) in \(R\).  If \(\theta^2=0\) and \(l=1+\theta\), then
\begin{equation}\label{eq:carry}
 N_4(l)=2\theta,
 \qquad
 l^4=1,
 \qquad
 N_8(l)=0.
\end{equation}
\end{lemma}

\begin{proof}
Since \(l^i=1+i\theta\),
\[
 N_4(l)=4+(1+2+3)\theta=4+6\theta=2\theta.
\]
Also \(l^4=1+4\theta=1\), and hence
\[
 N_8(l)=N_4(l)(1+l^4)=2N_4(l)=4\theta=0.
\]
\end{proof}

Thus a counterexample of the desired kind will follow from four pieces of
data:
\begin{equation}\label{eq:design-target}
 \boxed{
 \begin{gathered}
 4=0,\qquad \lambda=1+\theta,\qquad \theta^2=0,\\
 \Delta(U)=U\ot1+\lambda\ot U,
 \qquad 2\theta U\ne0.
 \end{gathered}}
\end{equation}
The next section explains why two quadratic coordinates are enough to
realize these requirements.

\section{Two rank-two cancellations and one bridge}

The elementary rank-two calculation is worth separating from the final
example.  Suppose that \(T^2=aT\), set \(g=1+bT\), and try to multiply two
copies of \(T\) by the rule
\[
 T_*=T_1+g_1T_2.
\]
Then
\begin{equation}\label{eq:rank-two-defect}
 T_*^2-aT_*
 =g_1\bigl(2T_1+a(g_1-1)\bigr)T_2
 =g_1(2+ab)T_1T_2.
\end{equation}
Consequently the single scalar identity
\begin{equation}\label{eq:rank-two-cancellation}
 ab=-2
\end{equation}
cancels the cross term.  Under this identity \(g\) is group-like,
\(g^2=1\), and
\[
 \Delta(T)=T\ot1+g\ot T
\]
is the familiar rank-two Hopf pattern.  Indeed,
\[
 \Delta(g)=g\ot g,
 \qquad
 [2]^\#(T)=(1+g)T=(2+ab)T=0,
\]
so coassociativity is formal and one may take \(S(T)=T\).  This is a
special case of the rank-two constructions appearing in the Oort--Tate
classification
\cite{OortTate}, but only the one-line calculation
\eqref{eq:rank-two-defect} is needed here.

To obtain rank four, begin formally with two such patterns.  Write
\[
 U^2=\alpha U+\beta V,
 \qquad
 V^2=\gamma V,
 \qquad
 \lambda=(1+rU)(1+sV).
\]
The term \(\beta V\) is the bridge between the two coordinates.  Expanding
the bridge ansatz suggests the following compact compatibility list:
\begin{equation}\label{eq:five-conditions}
\begin{array}{c@{\qquad}c@{\qquad}l}
 \alpha r=-2,&\gamma s=-2,
   &\text{the two diagonal cancellations},\\[2pt]
 \gamma r=0,&\beta s=0,
   &\text{the two mixed vanishings},\\[2pt]
 \multicolumn{2}{c}{\beta r=-\alpha s}
   &\text{the bridge balance}.
\end{array}
\end{equation}
The list separates the two diagonal cancellations, the two mixed
vanishings, and the bridge balance.  We only need the specialization below;
Section~\ref{sec:ambient} proves its sufficiency structurally through an
ambient matrix group, and Appendix~\ref{app:closure} records the complete
defect calculation.

There is a particularly economical way to solve
\eqref{eq:five-conditions}.  Take
\begin{equation}\label{eq:coefficient-specialization}
 \alpha=r^2,
 \qquad
 \beta=-rs,
 \qquad
 \gamma=s^2.
\end{equation}
The five conditions become
\[
 \underbrace{r^3=-2,\quad s^3=-2}_{\text{two cancellations}},
 \qquad
 \underbrace{rs^2=0}_{\text{both mixed vanishings}},
 \qquad
 -r^2s=-r^2s.
\]
The last balance is automatic.  This is the origin of the base
\[
 C=\Z[r,s]/(r^3+2,s^3+2,rs^2)
\]
and of the bridge equation
\[
 U^2=r^2U-rsV.
\]
No nilpotence cutoff has been added: the same three equations which make
the multiplication work will turn out to force the base to be Artinian.

The direct product of the two rank-two patterns would correspond to no
bridge.  That is not enough for our purpose.  The correction \(-rsV\) in
\(U^2\) changes the square of \(1+rU\) by the term which will eventually
become the nonzero fourth-power carry.  In this sense the bridge does not
merely couple the two blocks; it stores the obstruction.

\section{The base and its last infinitesimal layer}\label{sec:base}

We now study the three-relation base on its own.  Besides certifying the
nonvanishing needed later, its structure explains where the obstruction
lives.

\begin{proposition}\label{prop:base}
Let \(C\) be as in \eqref{eq:C-intro}.  Then \(C\) is a local Artin
complete intersection with maximal ideal
\[
 \m=(r,s)=(2,r,s)
\]
and residue field \(\F_2\).  As an abelian group,
\begin{equation}\label{eq:additive-C}
 C_{\mathrm{add}}
 \cong (\Z/4)\{1\}\oplus(\Z/4)\{s\}
 \oplus(\Z/2)\{s^2,r,rs,r^2,r^2s\}.
\end{equation}
In particular, \(C\) has characteristic four, cardinality \(2^9\), length
nine, and \(2s\ne0\).

The powers of the maximal ideal have successive \(\F_2\)-bases
\begin{equation}\label{eq:loewy-bases}
\begin{array}{c|c}
 \text{quotient}&\text{basis}\\ \hline
 C/\m&1\\
 \m/\m^2&r,\ s\\
 \m^2/\m^3&r^2,\ rs,\ s^2\\
 \m^3/\m^4&2,\ r^2s\\
 \m^4&2s.
\end{array}
\end{equation}
Thus \(\m^5=0\), the Hilbert function is \((1,2,3,2,1)\), and
\begin{equation}\label{eq:socle-C}
 \Soc(C)=\F_2(2s).
\end{equation}
More conceptually,
\begin{equation}\label{eq:gr-C}
 \gr_{\m}(C)
 \cong \F_2[R,S]/(RS^2,R^3+S^3).
\end{equation}
\end{proposition}

\begin{proof}
First put
\[
 C_0=\Z[r,s]/(r^3+2,s^3+2).
\]
Because both polynomials are monic, \(C_0\) is free over \(\Z\) on the
nine monomials \(r^is^j\), \(0\le i,j\le2\).  Let \(t=rs^2\).  On the
ordered basis
\[
 (1,s,s^2,r,rs,rs^2,r^2,r^2s,r^2s^2),
\]
multiplication by \(t\) gives
\[
 (rs^2,-2r,-2rs,r^2s^2,-2r^2,-2r^2s,-2s^2,4,4s).
\]
These nine entries lie in nine distinct basis directions.  Quotienting by
\((t)\) therefore gives exactly the additive normal form
\eqref{eq:additive-C}.  In particular \(2s\ne0\), while \(4=0\).

For later use, the three defining relations also give directly
\begin{equation}\label{eq:base-consequences}
 2r=0,
 \qquad
 2s^2=0,
 \qquad
 4=0.
\end{equation}
For example,
\[
 2r=r(s^3+2)-s(rs^2),
 \qquad
 2s^2=s^2(r^3+2)-r^2(rs^2),
\]
and \(4=0\) follows by multiplying \(2r=0\) by \(r^2\).

Since \(2=-r^3\), the ideals \((r,s)\) and \((2,r,s)\) agree.  Moreover
\(r^4=0\), \(s^5=0\), and \(C/(r,s)=\F_2\).  Thus \(C\) is local Artin with
the stated maximal ideal.  Multiplying the normal forms gives
\begin{align*}
 \m^2&=\F_2\{r^2,rs,s^2,r^2s,2,2s\},\\
 \m^3&=\F_2\{2,r^2s,2s\},\\
 \m^4&=\F_2\{2s\},
 \qquad \m^5=0,
\end{align*}
which proves \eqref{eq:loewy-bases}.

For completeness, the three defining equations form a regular sequence.
The first is a nonzerodivisor in \(\Z[r,s]\); modulo it, \(s^3+2\) is monic
in \(s\), hence is again a nonzerodivisor.  The resulting ring \(C_0\) is
free over \(\Z\), and \(rs^2\) becomes a unit after tensoring with
\(\mathbb Q\).  It is therefore a nonzerodivisor on \(C_0\).  Hence \(C\)
is a zero-dimensional complete intersection and is Gorenstein.  The
nonzero element \(2s\) is annihilated by \(r\) and \(s\), so the
one-dimensional socle is \eqref{eq:socle-C}.

Finally, the classes of \(r,s\) induce a surjection from the right side of
\eqref{eq:gr-C} to \(\gr_{\m}(C)\): the initial relations are \(RS^2=0\)
and \(R^3+S^3=0\).  The source is a homogeneous complete intersection of
two cubics, with Hilbert function \((1,2,3,2,1)\).  The equality of lengths
shows that the surjection is an isomorphism.
\end{proof}

The top class has a useful interpretation.  In the tangent cone it is
\(R^3S=S^4\); in \(C\) it lifts, up to sign, to
\[
 r^3s=-2s.
\]
Thus the two equations \(r^3=s^3=-2\) identify the two cubic directions
with the same arithmetic element, while the mixed equation \(rs^2=0\)
leaves their product with \(s\) as the unique final class.  This is exactly
the coefficient which appears in the fourth-power word.

\begin{corollary}[Quotient-minimality of this witness]\label{cor:quotient-minimal}
Every nonzero ideal of \(C\) contains \(2s\).  Consequently every proper
quotient \(C/J\) with \(0\ne J\subsetneq C\) kills the fourth-power defect
constructed below.  Equivalently, if a unital map \(C\to D\) sends \(2s\)
to a nonzero element, then the map is injective.
\end{corollary}

\begin{proof}
Every nonzero ideal in an Artin local ring meets the socle: multiply a
nonzero element by a largest possible power of the maximal ideal.  Now use
\eqref{eq:socle-C}.
\end{proof}

\section{The ambient matrix group}\label{sec:ambient}

We next explain conceptually why the coproduct in Theorem~\ref{thm:main}
is coassociative.  The point is not to check coassociativity coefficient by
coefficient.  Rather, the two coordinates form a closed submonoid of a
simple matrix group, so associativity is inherited.

We begin with a few identities in \(A\).  Monic reduction in \(V\) and
then in \(U\) gives
\begin{equation}\label{eq:A-free}
 A=C\{1,U,V,W\},\qquad W=UV.
\end{equation}
The basic products are
\begin{equation}\label{eq:W-products}
 UW=r^2W,
 \qquad
 VW=s^2W,
 \qquad
 W^2=0.
\end{equation}
Set
\begin{equation}\label{eq:L-M-theta}
 L=1+rU,
 \qquad
 M=1+sV,
 \qquad
 \lambda=LM=1+\theta.
\end{equation}

\begin{lemma}\label{lem:theta}
In \(A\) one has
\begin{equation}\label{eq:factor-squares}
 M^2=1,
 \qquad
 L^2=1+2sV,
 \qquad
 \lambda^2=1+2sV,
\end{equation}
and
\begin{equation}\label{eq:theta-core}
 2\theta=2sV,
 \qquad
 \theta^2=0,
 \qquad
 \lambda^{-1}=1-\theta.
\end{equation}
\end{lemma}

\begin{proof}
The identity \(V^2=s^2V\), together with \(s^3=-2\), gives
\[
 (1+sV)^2=1+(2s+s^4)V=1.
\]
Similarly, the bridge term in \(U^2=r^2U-rsV\) gives
\[
 (1+rU)^2=1+r^2U^2
 =1+r^4U-r^3sV=1+2sV,
\]
where \(2r=0\).  This proves \eqref{eq:factor-squares}.  Since
\[
 \theta=rU+sV+rsW,
\]
the base identities give \(2\theta=2sV\).  Comparing
\((1+\theta)^2\) with \(\lambda^2=1+2sV\) now gives \(\theta^2=0\), and the
last formula follows.
\end{proof}

The two factors \(L\) and \(M\) remember the two rank-two building blocks,
but after the bridge is inserted neither factor is itself a character.
Their product \(\lambda\) will be the character.  This distinction is the
reason for using an ambient group with one multiplicative coordinate.

Let
\[
 H=\Spec C[u,v,z,z^{-1}]
\]
with multiplication
\begin{equation}\label{eq:H-law}
 (u,v,z)(u',v',z')
 =\bigl(u+zu',\ vz'+v',\ zz'\bigr).
\end{equation}
This is visibly associative from the block-diagonal matrix realization
\begin{equation}\label{eq:matrix-realization}
 (u,v,z)\longmapsto
 \begin{pmatrix}z&u\\0&1\end{pmatrix}
 \oplus
 \begin{pmatrix}1&v\\0&z\end{pmatrix}.
\end{equation}
After replacing \(v\) by \(v/z\), it is the semidirect product
\(\Ga^2\rtimes\Gm\) in which \(\Gm\) acts with weights \(1\) and \(-1\).

Because \(\lambda\) is a unit, \(G_0=\Spec A\) is the closed subscheme of
\(H\) defined by
\begin{equation}\label{eq:G-in-H}
 u^2-r^2u+rsv=0,
 \qquad
 v^2-s^2v=0,
 \qquad
 z=(1+ru)(1+sv).
\end{equation}
At first \(G_0\) is only a pointed scheme.  The next proposition gives it
an associative multiplication.

\begin{proposition}\label{prop:submonoid}
The closed subscheme \(G_0\subset H\) contains the identity and is stable
under the multiplication of \(H\).  It is therefore a closed submonoid
scheme.  On its coordinate algebra, the induced coproduct is
\begin{equation}\label{eq:Delta-main}
 \Delta(U)=U\ot1+\lambda\ot U,
 \qquad
 \Delta(V)=V\ot\lambda+1\ot V,
 \qquad
 \Delta(\lambda)=\lambda\ot\lambda.
\end{equation}
\end{proposition}

\begin{proof}
The identity \((0,0,1)\) plainly satisfies \eqref{eq:G-in-H}.  For the
product of two points of \(G_0\), write
\begin{equation}\label{eq:product-coordinates}
 U_*=U_1+\lambda_1U_2,
 \qquad
 V_*=V_1\lambda_2+V_2.
\end{equation}
The two quadratic relations for \(U_*,V_*\) follow from the two rank-two
cancellations and the bridge identities; the exact two-line reduction is
given in Appendix~\ref{app:closure}.

The graph equation has a factorization which displays the role of the
bridge.  Put \(L_i=1+rU_i\), \(M_i=1+sV_i\), and
\[
 E=rsV_1U_2.
\]
The ambient product gives
\[
 1+rU_*=L_1(1+rM_1U_2),
 \qquad
 1+sV_*=(1+sV_1L_2)M_2.
\]
It follows that
\begin{equation}\label{eq:graph-factor}
 (1+rU_*)(1+sV_*)-\lambda_1\lambda_2
 =L_1M_2E(2+rU_2+sV_1+E).
\end{equation}
Every term on the right vanishes.  Indeed,
\[
 2E=E(sV_1)=E^2=0,
\]
using \(2r=rs^2=0\), while
\[
 E(rU_2)
 =r^2sV_1(r^2U_2-rsV_2)=0
\]
using \(2r=2s^2=0\).  Thus the graph equation is stable and
\(\lambda_*=\lambda_1\lambda_2\).  The multiplication and associativity
are inherited from \(H\), and the coordinate formulas are
\eqref{eq:Delta-main}.
\end{proof}

This proposition is the conceptual source of the bialgebra structure.  In
particular, coassociativity is no longer a separate identity to be checked:
it is simply associativity of the matrices in
\eqref{eq:matrix-realization}.  We have deliberately called \(G_0\) a
submonoid rather than a subgroup; inverses will emerge from the power
calculation itself.

\section{The fourth-power carry and the existence of inverses}
\label{sec:powers}

Proposition~\ref{prop:submonoid} and Lemma~\ref{lem:norm} give, already in
the monoid \(G_0\),
\begin{equation}\label{eq:norm-both}
 [n]^\#(U)=N_n(\lambda)U,
 \qquad
 [n]^\#(V)=VN_n(\lambda).
\end{equation}
By Lemmas~\ref{lem:theta} and \ref{lem:carry},
\begin{equation}\label{eq:N4-lambda}
 N_4(\lambda)=2\theta=2sV.
\end{equation}
Consequently
\begin{equation}\label{eq:fourth-U}
 [4]^\#(U)=2sVU=2sW.
\end{equation}
This is nonzero: \(A\) is free on \(1,U,V,W\), and Proposition
\ref{prop:base} proves \(2s\ne0\).

The other coordinate vanishes:
\begin{equation}\label{eq:fourth-V}
 [4]^\#(V)=2sV^2=2s^3V=-4V=0.
\end{equation}
Since the comorphism of a word is an algebra map,
\([4]^\#(W)=0\) as well.  This already proves that the fourth-power word is
not the unit word.

The same calculation at eight has an additional consequence.  Since
\(\lambda^4=1\),
\[
 N_8(\lambda)=N_4(\lambda)(1+\lambda^4)=0.
\]
Both \(U\) and \(V\) go to zero under \([8]^\#\), so
\begin{equation}\label{eq:eighth-unit}
 [8]=\unit
\end{equation}
as a word of the monoid \(G_0\).

\begin{proposition}\label{prop:inverse-by-seven}
The monoid scheme \(G_0\) is a group scheme.  Its inverse morphism is the
seventh-power word.
\end{proposition}

\begin{proof}
For every \(C\)-algebra \(R\) and every \(g\in G_0(R)\), associativity and
\eqref{eq:eighth-unit} give
\[
 g\,g^7=g^8=1=g^7g.
\]
Thus the morphism \([7]\colon G_0\to G_0\) is a functorial two-sided
inverse.  Hence \(G_0\) is a group object.
\end{proof}

We may therefore write \(G=G_0\).  On coordinate rings,
Proposition~\ref{prop:inverse-by-seven} supplies the antipode and completes
the Hopf algebra proof.  If an explicit formula is desired, the ambient
inverse gives
\begin{equation}\label{eq:explicit-antipode}
 S(U)=-\lambda^{-1}U=U-sW+r^2sV,
 \qquad
 S(V)=-V\lambda^{-1}=V+rW.
\end{equation}
These formulas are a check, not the reason inverses exist.

The nonzero class in \eqref{eq:fourth-U} is as deep as possible.

\begin{proposition}\label{prop:socle-A}
The ring \(A\) is a local Artin complete intersection of length \(36\), and
\begin{equation}\label{eq:socle-A}
 \Soc(A)=\F_2(2sW).
\end{equation}
In particular, the only nonzero coordinate of the fourth-power word spans
the socle of the total coordinate algebra.
\end{proposition}

\begin{proof}
The special fiber of \(A\) is local, so the finite algebra \(A\) over the
local ring \(C\) is local, with maximal ideal
\(\n=(r,s,U,V)\).  Its length is \(4\length(C)=36\).  Starting from the
regular sequence defining \(C\), first adjoin the monic relation
\(V^2-s^2V\), then the monic relation \(U^2-r^2U+rsV\).  Thus \(A\) is a
zero-dimensional complete intersection and has one-dimensional socle.

The element \(2sW\) is nonzero, and the identities
\[
 2rsW=2s^2W=2sUW=2sVW=0
\]
follow from \eqref{eq:base-consequences} and
\eqref{eq:W-products}.  It is therefore annihilated by \(\n\) and must
span the socle.
\end{proof}

We can now read the construction in one sentence.  The bridge turns the
failure of \(L=1+rU\) to remain an involution into the square-zero term
\(2sV\); the four-term geometric sum remembers that term, and
multiplication by \(U\) places it in the unique top class \(2sUV\).  The
eight-term sum sees the same carry twice and cancels it.

\section{A square-zero deformation and its obstruction class}
\label{sec:square-zero}

There are three useful scales at which to view the deformation.  The first
removes exactly the mixed bridge.  Put
\[
 q=rs,
 \qquad Q=(q)\subset C,
 \qquad D=C/Q.
\]
Since \(q^2=r^2s^2=r(rs^2)=0\), the map \(C\to D\) is a square-zero
extension.  Moreover \(2s=-r^2q\), so it already kills the fourth-power
class.

\begin{proposition}[The bridge-free matched product]
\label{prop:matched-product}
Over \(D\), the coordinate algebra becomes
\begin{equation}\label{eq:bridge-free-algebra}
 A_D
 \cong
 D[U,V]/(U^2-r^2U,\ V^2-s^2V),
\end{equation}
and \(L=1+rU\) and \(M=1+sV\) are separate group-like involutions.  The
closed rank-two subgroup schemes
\[
 H_U=(V=0),
 \qquad H_V=(U=0)
\]
are both killed by two, and multiplication induces an isomorphism of
\(D\)-schemes
\begin{equation}\label{eq:matched-factorization}
 H_U\mathbin{\times}_D H_V\xrightarrow{\ \sim\ }G_D.
\end{equation}
Thus \(G_D\) is the exact matched product of its two elementary rank-two
factors.

Upstairs, the failures of the two factors to remain group-like are the
\(Q\)-valued terms
\begin{equation}\label{eq:character-cocycles}
 \begin{split}
  c_L&=\Delta(L)-L\ot L=q\,LV\ot U,\\
  c_M&=\Delta(M)-M\ot M=q\,V\ot UM.
 \end{split}
\end{equation}
They are normalized twisted co-Hochschild \(2\)-cocycles, and their weighted
sum vanishes so that \(LM=\lambda\) remains group-like.
\end{proposition}

\begin{proof}
Setting \(q=0\) removes the bridge from the relation for \(U^2\), proving
\eqref{eq:bridge-free-algebra}.  The formulas of Lemma~\ref{lem:theta} give
\(L^2=M^2=1\), because \(2s=-r^2q=0\) in \(D\).  For group-likeness, the
product coordinates give
\begin{align*}
 \Delta(L)-L_1L_2
 &=rL_1(M_1-1)U_2=qL_1V_1U_2,\\
 \Delta(M)-M_1M_2
 &=sV_1(L_2-1)M_2=qV_1U_2M_2.
\end{align*}
These are both zero over \(D\), and over \(C\) they are exactly
\eqref{eq:character-cocycles}.

The two quotient Hopf algebras are the familiar rank-two algebras
\(D[U]/(U^2-r^2U)\) and \(D[V]/(V^2-s^2V)\).  The rank-two cancellation
\((2+rU)U=(2+r^3)U=0\), and its \((s,V)\) analogue, shows that both subgroup
schemes are killed by two.  As an algebra, \(A_D\) is the tensor product of
these two rank-two coordinate algebras.  More geometrically, the ambient
law gives the unique factorization
\[
 (u,v,LM)=(u,0,L)(0,v,M),
\]
which proves \eqref{eq:matched-factorization}.  Reversing the order produces
the mutual actions, so this is a matched product rather than a claimed
direct product of groups.

Finally, coassociativity of \(\Delta\) and \(Q^2=0\) linearize to
\[
 (\overline\Delta\ot\id)c_L+c_L\ot\overline L
 = (\id\ot\overline\Delta)c_L+\overline L\ot c_L,
\]
and similarly for \(c_M\); the counit makes them normalized.  This is the
twisted co-Hochschild cocycle equation.  Expanding
\(\Delta(LM)=\Delta(L)\Delta(M)\) to first order in \(Q\), and using
\(\Delta(\lambda)=\lambda\ot\lambda\), gives
\[
 c_L(\overline M\ot\overline M)
 +(\overline L\ot\overline L)c_M=0.
\]
Thus the two individual character defects cancel in their product.
\end{proof}

\subsection{What is being deformed?}
\label{subsec:what-deforms}

There is an important distinction hidden in
\eqref{eq:matched-factorization}.  The multiplication map identifies
\(G_D\) with \(H_U\times_DH_V\) as a \emph{scheme}, but not in general as
a group scheme.  The group law remembers mutual actions between the two
factors; equivalently, \(G_D\) is an exact matched product.  Only after
passing all the way to the closed fiber do the actions disappear and give
the literal direct product
\(\alpha _2\times\alpha _2\).  Thus the extension
\[
 0\longrightarrow Q\longrightarrow C\longrightarrow D\longrightarrow0
\]
asks us to lift an exact factorization, not merely to lift a direct
product.

Three kinds of data can change in such a lift.  The two rank-two group
schemes can deform separately; their mutual actions, and hence the
coproduct on their tensor-product coordinate algebra, can deform; and the
tensor-product algebra itself can acquire a mixed multiplication term.
The last kind of term is the bridge.  These are useful descriptions, but
they are not three independent summands of a deformation space.  The
identity that the coproduct is an algebra map couples multiplication and
coproduct corrections.  A bridge in the algebra will usually force a
mixed correction to the coproduct.

Here is the standard cohomological formulation.  Put \(A_0=A_D\) and
\(M=Q\ot_DA_0\simeq QA\).  Frame a lift by identifying its underlying
\(C\)-module with a chosen lift of the free \(D\)-module \(A_0\).  If
\((m_1,\Delta_1)\) and \((m_2,\Delta_2)\) are two \emph{actual} framed
bialgebra lifts, their differences descend to \(D\)-linear maps
\[
 \mu\colon A_0\ot_DA_0\longrightarrow M,
 \qquad
 \delta\colon A_0\longrightarrow
 Q\ot_D(A_0\ot_DA_0).
\]
The hypothesis \(Q^2=0\) makes every bialgebra identity linear in
\((\mu,\delta)\).  Thus \(\mu\) is a normalized Hochschild \(2\)-cocycle,
\(\delta\) is a normalized co-Hochschild \(2\)-cocycle, and they satisfy
the mixed equation
\begin{align}
 \delta(ab)+\overline\Delta\bigl(\mu(a,b)\bigr)
 &=
 \delta(a)\overline\Delta(b)
 +\overline\Delta(a)\delta(b)
 +\mu^{(2)}(\overline\Delta(a),\overline\Delta(b)),
 \label{eq:GS-mixed}
\\[-2pt]
 \mu^{(2)}(a_1\ot a_2,b_1\ot b_2)
 &=
 \mu(a_1,b_1)\ot a_2b_2
 +a_1b_1\ot\mu(a_2,b_2).
 \nonumber
\end{align}
For completeness, the coalgebra equation is
\begin{equation}\label{eq:GS-coalgebra}
 (\overline\Delta\ot\id)\delta
 +(\delta\ot\id)\overline\Delta
 =
 (\id\ot\overline\Delta)\delta
 +(\id\ot\delta)\overline\Delta .
\end{equation}
A framed coordinate change \(\id+h\), where
\(h\colon A_0\to M\) is normalized, changes the pair by
\begin{align}
 \mu(a,b)&\longmapsto
 \mu(a,b)+a h(b)-h(ab)+h(a)b,\label{eq:GS-gauge-m}\\
 \delta&\longmapsto
 \delta+\overline\Delta h
 -(h\ot\id+\id\ot h)\overline\Delta .
 \label{eq:GS-gauge-d}
\end{align}
These are the Gerstenhaber--Schack cocycle and coboundary equations
\cite{GerstenhaberSchack}.  Because an affine group scheme has a
commutative coordinate algebra, our deformation problem is the
commutative subproblem in which \(\mu\) is symmetric.  In this language
infinitesimal automorphisms, deformations, and existence obstructions are
measured in degrees \(1,2,3\), respectively.

There are two qualifications to this slogan.  First, \(C\to D\) is not a
split square-zero extension.  A coefficientwise lift of the formulas over
\(D\) need not itself satisfy the bialgebra identities, so its correction
obeys an \emph{inhomogeneous} version of
\eqref{eq:GS-mixed}.  Once one actual lift has been chosen, differences
from it are honest cocycles.  The set of lifts is therefore naturally
affine, or torsorial, rather than a vector space with a canonical zero.
Second, no separate infinitesimal antipode has to be chosen.  In the
convolution algebra, an inverse modulo a nilpotent ideal lifts by a finite
geometric series, so a bialgebra lift of a Hopf algebra is again Hopf.

For our chosen coefficientwise lift of the bridge-free presentation, the
visible corrections are
\begin{equation}\label{eq:visible-GS-corrections}
 \mu(U,U)=-qV,
 \qquad
 \delta(U)=qW\ot U,
 \qquad
 \delta(V)=qV\ot W.
\end{equation}
Here the last two terms come from
\(\lambda=1+rU+sV+qW\).  They are components of one compatible
bialgebra correction, not three unrelated choices.  The character
cocycles \(c_L,c_M\) of \eqref{eq:character-cocycles} give a more
intrinsic view of its coalgebra part: neither of the two rank-two
characters remains group-like, but their defects cancel so that
\(LM=\lambda\) does.

\subsection{An explicit chart of possible lifts}
\label{subsec:explicit-lifts}

The general cohomology is useful for orientation, but in the present
example the relevant part of the deformation problem can be written as a
five-term linear calculation.  Consider the triangular bridge chart
\[
 U^2=\alpha U+\beta V,\qquad V^2=\gamma V,\qquad
 \lambda=(1+rU)(1+sV),
\]
with the same skew-primitive formulas for \(U,V\).  The five compatibility
equations are those of \eqref{eq:five-conditions}:
\begin{equation}\label{eq:bridge-chart-again}
 \alpha r=-2,\quad \gamma s=-2,\quad
 \gamma r=0,\quad \beta s=0,\quad
 \beta r+\alpha s=0.
\end{equation}

\begin{lemma}[The bridge chart is Hopf]\label{lem:bridge-chart-Hopf}
The five equations \eqref{eq:bridge-chart-again} make both quadrics stable
under the coproduct and make \(\lambda\) group-like.  They therefore
define a rank-four Hopf algebra.  Its antipode is
\begin{equation}\label{eq:general-bridge-antipode}
 S(U)=U-sUV+\alpha sV,\qquad S(V)=V+rUV.
\end{equation}
\end{lemma}

\begin{proof}
We give the calculation because it displays how the bridge equation enters
the Hopf axioms.  The five scalar identities first imply
\[
 2r=4=\alpha s^2=2s^2=0,
 \qquad
 r^2\beta=2s.
\]
Put \(L=1+rU\), \(M=1+sV\), and \(\theta=\lambda-1\).  Then
\begin{equation}\label{eq:general-bridge-identities}
 \begin{gathered}
  M^2=1,\qquad L^2=\lambda^2=1+2sV,\qquad \theta^2=0,\\
  \gamma(\lambda-1)=-2V,\qquad
  \beta(\lambda-1)=-\alpha sU,\qquad
  \beta(\lambda^2-1)=0,\\
  2\lambda U+\alpha(\lambda^2-\lambda)=-\alpha sV.
 \end{gathered}
\end{equation}
For example, the last identity follows by adding
\[
 2\lambda U=2U-2\alpha sV+2sUV
\]
and
\[
 \alpha(\lambda^2-\lambda)=2U+\alpha sV+2sUV.
\]

On two universal points write
\[
 U_*=U_1+\lambda_1U_2,\qquad
 V_*=V_1\lambda_2+V_2.
\]
Reduction by the two monic quadrics gives
\begin{align*}
 V_*^2-\gamma V_*
 &=V_1\lambda_2
   \bigl(\gamma(\lambda_2-1)+2V_2\bigr),\\
 U_*^2-\alpha U_*-\beta V_*
 &=\beta V_1(1-\lambda_2)\\
 &\quad+
 \bigl(2\lambda_1U_1+\alpha(\lambda_1^2-\lambda_1)\bigr)U_2
 +\beta(\lambda_1^2-1)V_2.
\end{align*}
The first line vanishes by \eqref{eq:general-bridge-identities}.  In the
second, the first two terms are respectively
\(\alpha sV_1U_2\) and \(-\alpha sV_1U_2\), while the last is zero.
Thus both relations are preserved.

It remains to check the graph coordinate.  Put \(t=rs\).  Directly from
the definitions,
\[
 \Delta(L)=L_1L_2+tL_1V_1U_2,\qquad
 \Delta(M)=M_1M_2+tV_1U_2M_2.
\]
After factoring out \(L_1M_2\), the error in
\(\Delta(LM)=L_1M_1L_2M_2\) is
\[
 tV_1U_2(L_2+M_1)+t^2V_1^2U_2^2.
\]
It vanishes after substituting \(L_2+M_1=2+rU_2+sV_1\):
the four resulting coefficients reduce respectively by
\[
 2t=0,\qquad tr\alpha=tr\beta=0,\qquad ts\gamma=0,
 \qquad t\gamma=0.
\]
Hence \(\lambda=LM\) is group-like.  Coassociativity is now formal from
the two skew-primitive formulas, and the counit is the evident one.

Finally, \(4=0\) and \(\theta^2=0\) imply
\(\lambda^4=1\) and \(N_8(\lambda)=0\) by
Lemma~\ref{lem:carry}.  The skew-primitive power formulas show that the
eighth-power word is the unit word.  The seventh-power word is therefore
an inverse and supplies the antipode.  Reducing its two coordinates gives
\eqref{eq:general-bridge-antipode}.
\end{proof}

This chart does not claim to classify every abstract rank-four Hopf
deformation; it classifies the deformations which preserve the two named
axes, the triangular quadratic normal form, and the skew-primitive
coproduct.

It is helpful to record its tangent complex.  Let \(R'\to R\) be a
square-zero extension with kernel \(J\), and start at a bridge-free point
\[
 (\overline\alpha,0,\overline\gamma,\overline r,\overline s)
\]
of \eqref{eq:bridge-chart-again}.  The difference between two lifts is a
tuple
\[
 (a,b,c,\rho,\sigma)\in J^5,
\]
corresponding in order to
\((\alpha,\beta,\gamma,r,s)\).  Linearizing the five equations gives
\begin{equation}\label{eq:explicit-deformation-d1}
 d_1
 \begin{pmatrix}a\\ b\\ c\\ \rho\\ \sigma\end{pmatrix}
 =
 \begin{pmatrix}
  \overline r a+\overline\alpha\rho\\
  \overline s c+\overline\gamma\sigma\\
  \overline r c+\overline\gamma\rho\\
  \overline s b\\
  \overline s a+\overline r b+\overline\alpha\sigma
 \end{pmatrix}.
\end{equation}
The diagonal coordinate changes
\[
 U'=xU,\qquad V'=yV
\]
act by
\[
 (\alpha,\beta,\gamma,r,s)\longmapsto
 (x\alpha,x^2y^{-1}\beta,y\gamma,x^{-1}r,y^{-1}s).
\]
Writing \(x=1+\xi\) and \(y=1+\eta\), their infinitesimal action at the
bridge-free point is
\begin{equation}\label{eq:explicit-deformation-d0}
 d_0(\xi,\eta)=
 (\overline\alpha\xi,0,\overline\gamma\eta,
  -\overline r\xi,-\overline s\eta).
\end{equation}
Thus
\[
 J^2\xrightarrow{\ d_0\ }J^5\xrightarrow{\ d_1\ }J^5
\]
is a concrete presentation-level shadow of the
Gerstenhaber--Schack complex.  If one lift exists, its framed companions
form an affine space under \(\ker(d_1)\), and quotienting by
\(\operatorname{im}(d_0)\) accounts for these coordinate gauges.
For a nonsplit base extension, a naive coefficientwise lift has a
curvature vector \(o\in J^5\), and the actual lifting equation is
\[
 d_1(a,b,c,\rho,\sigma)=-o.
\]

In our case the downstairs point is
\[
 (r^2,0,s^2,r,s)\quad\text{over }D.
\]
Its naive lift to \(C\), still with \(\beta=0\), violates exactly the last
equation of \eqref{eq:bridge-chart-again}:
\[
 \beta r+\alpha s=r^2s=rq.
\]
At the level of the coproduct the same curvature is the single defect
\begin{equation}\label{eq:naive-lift-curvature}
 \Delta(U)^2-r^2\Delta(U)=rq\,V\ot U.
\end{equation}
The other quadratic relation and the graph equation for \(\lambda\)
remain valid.  Thus \(rq\) is the primary mixed obstruction to retaining
the tensor-product algebra upstairs.  The correction
\[
 (a,b,c,\rho,\sigma)=(0,-q,0,0,0)
\]
has \(d_1(0,-q,0,0,0)=(0,0,0,0,-rq)\), so the bridge is literally a
cochain which null-homotopes this curvature.

\begin{proposition}[The triangular lifts across \(Q\)]
\label{prop:triangular-lifts}
Keep \(r,s\) fixed and consider all lifts in the triangular chart.  They
are precisely
\begin{equation}\label{eq:eight-triangular-lifts}
 \alpha=r^2+a,\qquad
 \gamma=s^2+c,\qquad
 \beta=-q+j,
 \qquad
 a,c,j\in(2s).
\end{equation}
Since \((2s)\simeq\F_2\), these are eight framed normal-form lifts.  If the
two diagonal rank-two data are also fixed, there are exactly two:
\[
 \beta=-q\quad\text{or}\quad \beta=-q+2s.
\]
No assertion is made here that the eight framings give eight abstract
isomorphism classes.

Every lift in \eqref{eq:eight-triangular-lifts} has the same nonzero
fourth-power word:
\begin{equation}\label{eq:all-lifts-same-carry}
 [4]^\#(U)=2sW,\qquad [4]^\#(V)=0.
\end{equation}
\end{proposition}

\begin{proof}
As a \(D\)-module, the square-zero kernel is
\begin{equation}\label{eq:Q-module-chain}
 Q=\F_2\{q,rq,r^2q\},
 \qquad
 q\xmapsto{\ r\ }rq
 \xmapsto{\ r\ }r^2q=-2s
 \xmapsto{\ r\ }0,
\end{equation}
and \(sQ=2Q=0\).  Hence
\[
 \Ann_Q(r)=(r^2q)=(2s).
\]
This also proves that \(C\to D\) is nonsplit.  Indeed, any lifts of
\(\overline r,\overline s\) have the form \(r+a,s+b\) with \(a,b\in Q\),
and their product is \(q+rb\), because \(sQ=Q^2=0\).  It cannot vanish:
\(q\notin rQ=\F_2\{rq,r^2q\}\).
Write
\(\alpha=r^2+a\), \(\gamma=s^2+c\), and \(\beta\in Q\).
Substitution into \eqref{eq:bridge-chart-again} gives
\[
 ra=0,\qquad sc=0,\qquad rc=0,\qquad
 s\beta=0,\qquad r(\beta+q)+sa=0.
\]
The equations containing \(sQ\) are automatic, leaving
\[
 a,c\in\Ann_Q(r),
 \qquad
 \beta=-q+j,\quad j\in\Ann_Q(r).
\]
This proves \eqref{eq:eight-triangular-lifts}.

The bridge balance itself already contains the secondary invariant.
Multiplying \(\beta r=-\alpha s\) by \(r\) and using
\(\alpha r=-2\) gives
\begin{equation}\label{eq:secondary-bridge-residue}
 r^2\beta=-r\alpha s=2s.
\end{equation}
The two diagonal cancellations imply
\[
 (1+sV)^2=1,\qquad
 (1+rU)^2=1+r^2\beta V=1+2sV.
\]
Therefore \(\lambda^2=1+2sV\), while
\((\lambda-1)^2=0\), and Lemma~\ref{lem:carry} gives
\[
 N_4(\lambda)=2sV.
\]
The skew-primitive power formulas now yield
\[
 [4]^\#(U)=N_4(\lambda)U=2sW,
 \qquad
 [4]^\#(V)=V N_4(\lambda)=2s\gamma V=0.
\]
The first element is nonzero by freeness of \(A\) over \(C\) and
\(2s\ne0\).
\end{proof}

The proposition gives a particularly elementary cohomological picture.
If the two diagonal factors are fixed, a bridge is an element
\(\beta\in Q\) satisfying
\begin{equation}\label{eq:bridge-nullhomotopy}
 s\beta=0,\qquad r\beta=-rq.
\end{equation}
Thus the obstruction to combining the lifted rank-two blocks is the class
of \((0,-rq)\) in the cokernel of
\[
 Q\longrightarrow Q\oplus Q,\qquad
 b\longmapsto(sb,rb).
\]
When that class vanishes, the possible null-homotopies form a torsor under
\(\Ann_Q(r,s)\).  In our example this torsor has the two elements displayed
above.  Its secondary residue \(r^2\beta=2s\) is independent of the
null-homotopy, since two choices differ by an element killed by \(r\).
The power word identifies this well-defined secondary residue, multiplied
by \(W\), with the failure of killedness.

This behavior would be impossible for a split, factor-fixed perturbation
in which the downstairs equality \(\alpha s=0\) lifted unchanged.  The
bridge equations would then give \(r\beta=s\beta=0\), hence
\(r^2\beta=0\).  What makes the present construction work is the
nonsplit multiplication in the base: \(rs=0\) downstairs lifts to
\(rs=q\ne0\) upstairs.  Moreover \(Q\) is square-zero but not a small
one-step module; the three-term \(r\)-chain
\eqref{eq:Q-module-chain} remembers one further multiplication.  The
bridge kills the primary defect \(rq\), while that extra multiplication
leaves the secondary class \(2s\).

\begin{remark}[The exact factorization does not lift]
\label{rem:factorization-does-not-lift}
The factor \(H_U\) lifts, since \((V)\) remains a Hopf ideal.  There is,
however, no finite flat rank-two closed subgroup of \(G\) lifting the
complement \(H_V\).  Indeed, suppose its coordinate algebra \(E\) existed.
If \(u,v\) are the images of \(U,V\), then \(1,v\) is a \(C\)-basis after
Nakayama's lemma, and \(u\) reduces to zero modulo \(Q\).  We may therefore
write
\[
 u=q(a+bv).
\]
Since \(Q^2=0\), one has \(u^2=0\).  The relation
\(U^2-r^2U+qV=0\) then gives \(qv=r^2u\).  Comparing the coefficients of
\(v\) in the basis \(1,v\) gives
\[
 q(1-r^2b)=0.
\]
The factor \(1-r^2b\) is a unit, whereas \(q\ne0\) remains nonzero in the
free \(C\)-module \(E\), a contradiction.  Thus the multiplication
cocycle in \eqref{eq:visible-GS-corrections} measures the actual loss of
the exact factorization, not merely a change of coordinates.
\end{remark}

\subsection{Why this is a counterexample to the rank-times-power conjecture}
\label{subsec:why-counterexample}

The Hopf algebra and the equation \([4]=\unit\) have distinct deformation
theories.  The Gerstenhaber--Schack equations ask whether multiplication
and coproduct lift at all.  The bridge solves that problem.  Once a Hopf
lift exists, the additional equality of the fourth-power word with the
unit word has its own obstruction.

Indeed, \(G_D\) is killed by four, so the two algebra maps
\([4]^\#\) and \(e^\#=\eta\eps\) agree modulo \(Q\).  Their difference
\[
 \omega_4=[4]^\#-e^\#\colon A\longrightarrow QA
\]
kills \(QA\) and descends to an \(e\)-derivation on \(A_0\).  Explicitly,
\begin{equation}\label{eq:Q-power-obstruction}
 [\omega_4]\in
 \Der_D(A_0,(QA)_e)
 \cong
 \Hom_{A_0}\bigl(\Omega^1_{A_0/D},(QA)_e\bigr),
 \qquad
 \omega_4(U)=2sW,\quad \omega_4(V)=0.
\end{equation}
The derivation identity follows directly by expanding the difference of
two algebra maps; its quadratic error lies in \(Q^2A=0\).  The class is
nonzero.  After choosing one Hopf lift as origin, formation of the power
word therefore gives a linear infinitesimal operation from differences of
Hopf deformation classes to augmentation derivations.  Before choosing an
origin, it is more accurately an affine obstruction map on the torsor of
lifts.  Proposition~\ref{prop:triangular-lifts} says that its value is the
same nonzero class on every point of our explicit chart.

Grothendieck's rank-times-power question predicts that every finite
locally free group scheme of rank \(n\) should have trivial \(n\)-th power
word.  Here the two monic quadrics keep \(A\) finite free of rank four,
and all of the Hopf axioms hold, but
\[
 [4]^\#(U)=2sUV\ne0=e^\#(U).
\]
Thus \([4]\ne\unit\), which is exactly a counterexample in rank four.
Every geometric fiber kills \(Q\), and so sees none of this vertical
derivation.  Deligne's commutative theorem does not apply because the
coproduct deformation is noncocommutative.  In short, nilpotence allows
the nonzero vertical obstruction to survive, while noncommutativity allows
it to occur as a power word.

The mechanism may be summarized as
\[
 \boxed{
 \begin{gathered}
 \text{the lifted matched product has curvature }rq,\\
 r\beta=-rq\text{ makes the Hopf lift exist},\\
 r^2\beta=2s\text{ is the secondary residue},\\
 [4]^\#(U)=2sUV\ne0.
 \end{gathered}}
\]

At the first infinitesimal neighborhood one can also see the primary
bialgebra deformation datum directly.  Set
\[
 C^{(1)}=C/\m^2\cong\F_2[r,s]/(r,s)^2.
\]
There the underlying algebra is simply
\[
 A^{(1)}=C^{(1)}[U,V]/(U^2,V^2),
\]
and, if \(\Delta_0\) denotes the primitive coproduct, then
\begin{equation}\label{eq:first-order-coproduct}
 \Delta=\Delta_0+r\Phi_r+s\Phi_s,
 \qquad
 \begin{array}{c|cc}
  &U&V\\ \hline
  \Phi_r&U\ot U&V\ot U\\
  \Phi_s&V\ot U&V\ot V.
 \end{array}
\end{equation}
The requirement that \(\Delta\) be an algebra map and the coassociativity
identity, linearized modulo \((r,s)^2\), say exactly that
\(\Phi_r,\Phi_s\) are the coproduct components of first-order bialgebra
\(2\)-cocycles.  In Gerstenhaber--Schack terminology this is the primary
deformation datum \cite{GerstenhaberSchack}.  It is already
noncocommutative, while
\begin{equation}\label{eq:first-order-doubling}
 [2]^\#(U)=sW,
 \qquad
 [2]^\#(V)=rW,
 \qquad
 [4]=\unit.
\end{equation}
The bridge \(-rsV\) is invisible at this order: it is a second-order
compatibility correction.  Thus noncommutativity occurs in the primary
deformation, whereas failure of the fourth-power relation is a secondary
effect supported in the final Loewy layer.

A different square-zero quotient makes the power calculation itself
transparent.  Let \(\m=(r,s)\subset C\), and put
\begin{equation}\label{eq:square-zero-extension}
 I=\m^3=(2,r^2s,2s),
 \qquad
 B=C/I.
\end{equation}
The Loewy calculation in Proposition~\ref{prop:base} gives
\begin{equation}\label{eq:easy-base}
 I^2=0,
 \qquad
 B\cong \F_2[r,s]/(r,s)^3.
\end{equation}
Thus \(C\) is a square-zero thickening of a particularly transparent
characteristic-two base.  Write
\[
 A_0=A/IA,
 \qquad
 G_0=G\mathbin{\times}_{\Spec C}\Spec B,
 \qquad
 \phi=[2]\colon G\longrightarrow G.
\]

\begin{proposition}[The square-zero carry]\label{prop:square-zero-carry}
On \(A\), the doubling word is
\begin{equation}\label{eq:doubling-upstairs}
 \phi^\#(U)=-sW-r^2sV,
 \qquad
 \phi^\#(V)=rW,
 \qquad
 \phi^\#(W)=0.
\end{equation}
After reduction to \(A_0\), these formulas become
\begin{equation}\label{eq:doubling-downstairs}
 \overline\phi^{\,\#}(U)=sW,
 \qquad
 \overline\phi^{\,\#}(V)=rW,
 \qquad
 \overline\phi^{\,\#}(W)=0.
\end{equation}
Consequently \([4]=\unit\) on \(G_0\), whereas on its square-zero lift
\(G\) one has
\begin{equation}\label{eq:bockstein-carry}
 [4]^\#(U)
 =-r^2s\,\phi^\#(V)
 =-r^3sW
 =2sW\ne0.
\end{equation}
\end{proposition}

\begin{proof}
The equality \(I=\m^3\) is \eqref{eq:loewy-bases}; since \(\m^5=0\), it
implies \(I^2=0\).  Killing \(I\) kills \(2\) and all monomials of degree
three in \(r,s\), which proves the presentation of \(B\) in
\eqref{eq:easy-base}.

Since \(U\) and \(V\) are respectively left and right
\(\lambda\)-skew primitive,
\[
 \phi^\#(U)=(1+\lambda)U,
 \qquad
 \phi^\#(V)=V(1+\lambda).
\]
Using \(U^2=r^2U-rsV\), \(UW=r^2W\), and \(r^3=-2\), the first expression
reduces to
\begin{align*}
 (2+rU+sV+rsW)U
 &=2U+r^3U-r^2sV+sW+r^3sW\\
 &=-sW-r^2sV.
\end{align*}
Similarly, \(V^2=s^2V\), \(VW=s^2W\), \(s^3=-2\), and \(2r=0\) give
\(\phi^\#(V)=rW\).  Their product is zero, so
\(\phi^\#(W)=0\).  Modulo \(I\), the element \(r^2s\) and the scalar \(2\)
vanish, and minus equals plus; this proves
\eqref{eq:doubling-downstairs}.  In particular the image of doubling is
contained in the square-zero \(W\)-direction, which doubling itself kills.
Thus \(\overline\phi^2=\unit\).  Applying \(\phi^\#\) once more to the
first formula in \eqref{eq:doubling-upstairs} gives
\[
 \phi^{\#2}(U)
 =-s\phi^\#(W)-r^2s\phi^\#(V)
 =-r^3sW=2sW,
\]
and proves \eqref{eq:bockstein-carry}.
\end{proof}

The last line has the shape of a Bockstein calculation.  Downstairs,
doubling sends \(V\) to \(rW\).  The lift of the other doubling coordinate
contains the correction \(-r^2sV\in IA\).  Composing the two multiplies
that correction by \(r\), and the lifted cubic relation
\(r^3=-2\) turns it into the final class \(2sW\).  The square-zero ideal
\(I\) is not annihilated by \(\m\): indeed
\(r(r^2s)=-2s\ne0\).  This is exactly what permits a first-order
correction to fall one Loewy step farther and remain visible.

For the sharp cohomological statement, pass to the final Loewy layer.  Put
\begin{equation}\label{eq:minimal-square-zero}
 J=(c)=\m^4\subset C,
 \qquad c=2s,
 \qquad \overline C=C/J,
 \qquad \overline A=A/JA.
\end{equation}
Then \(J^2=0\), and in fact \(\m J=0\).  Let
\(e^\#\colon A\to A\) be the algebra map obtained by applying the counit
and then including \(C\) into \(A\).  The precise infinitesimal difference
between the two word maps is
\begin{equation}\label{eq:obstruction-derivation}
 D=[4]^\#-e^\#\colon A\longrightarrow JA.
\end{equation}
The target is correct because \([4]=\unit\) after reduction to
\(\overline C\).  Moreover \((JA)^2=0\), so the difference of the two
algebra maps satisfies
\begin{equation}\label{eq:e-derivation}
 D(ab)=e^\#(a)D(b)+D(a)e^\#(b).
\end{equation}
It kills \(JA\), and hence descends to an \(e\)-derivation on
\(\overline A\).  If \(M=JA\) is regarded as an \(\overline A\)-module
through the identity section, then the resulting canonical class is
\begin{equation}\label{eq:cotangent-class}
 [D]\in
 \operatorname{Der}_{\overline C}(\overline A,M_e)
 \cong
 \operatorname{Hom}_{\overline A}
       (\Omega^1_{\overline A/\overline C},M_e)
 \cong
 \mathrm H^0\operatorname{RHom}_{\overline A}
       (\mathbb L_{\overline A/\overline C},M_e).
\end{equation}
This is the standard cotangent-complex description of infinitesimal
differences between maps \cite{IllusieCotangent}.
In coordinates it is simply
\begin{equation}\label{eq:class-coordinates}
 D(U)=2sW,
 \qquad
 D(V)=0.
\end{equation}
For example, pulling the relative cotangent module back to the identity
gives
\[
 e^*\Omega^1_{\overline A/\overline C}
 \cong
 \frac{\overline C\,dU\oplus\overline C\,dV}
      {(r^2dU-rsdV,\ s^2dV)}.
\]
The assignment \(dU\mapsto2sW\), \(dV\mapsto0\) respects these relations,
since \(2r^2s=0\), and it is nonzero by
Proposition~\ref{prop:socle-A}.  Thus \([D]\) is the exact infinitesimal
obstruction to lifting the equality \([4]=\unit\) from
\(G_{\overline C}\) to \(G\).

This also locates the relevant cohomological degree.  The Hopf algebra
itself has already been lifted, so there is no obstruction to the existence
of the multiplication or antipode waiting in a higher deformation group.
What fails to lift is an equality between two morphisms.  Across a
square-zero extension, the difference between such morphisms is measured
by a derivation, hence by the degree-zero cotangent cohomology in
\eqref{eq:cotangent-class}.

The same minimal layer has an additional character interpretation.  The
base change of \(G\) to \(\overline C\) is killed by four.  On the lift,
there is first a character-valued obstruction:
\begin{equation}\label{eq:infinitesimal-character}
 \kappa(\lambda)=\lambda^2-1=cV.
\end{equation}
It is primitive.  Indeed \(c\theta=0\), hence \(c\lambda=c\), and
\begin{equation}\label{eq:cV-primitive}
 \Delta(cV)
 =c(V\ot\lambda+1\ot V)
 =cV\ot1+1\ot cV.
\end{equation}
For a square-zero ideal, the kernel of reduction on multiplicative units is
the additive group \(1+JA\).  Thus the nonzero primitive element \(cV\)
is precisely the infinitesimal character measuring the failure of the
order-two character \(\overline\lambda\) to lift as an order-two character.
This obstruction is independent of the chosen character lift: another
lift has the form \(\lambda(1+y)\), with \(y\in JA\) primitive, and
\((1+y)^2=1\) because \(2JA=(JA)^2=0\).  Hence every character lift has
square \(1+cV\).  Equivalently, \(\lambda\) factors through \(\mu_4\) but
not through \(\mu_2\).

The Bockstein terminology is literal.  Multiplication by two in
\(0\to JA\to A\to\overline A\to0\) gives
\begin{equation}\label{eq:literal-bockstein}
 \beta\colon \overline A[2]\longrightarrow JA,
 \qquad
 \beta(\overline x)=2x,
\end{equation}
where independence of the lift again uses \(2JA=0\).  For
\(\theta=\lambda-1\), one has \(2\overline\theta=0\) and
\begin{equation}\label{eq:theta-bockstein}
 \beta(\overline\theta)=2\theta=cV.
\end{equation}
The geometric sum converts this character obstruction into the power-word
obstruction:
\begin{equation}\label{eq:character-to-word-obstruction}
 N_4(\lambda)=2\theta=\beta(\overline\theta)=cV,
 \qquad
 D(U)=\kappa(\lambda)U=cW.
\end{equation}
This is the same cotangent class \((2sW,0)\) found above.

Since every nonzero ideal of \(C\) contains \(J\),
Corollary~\ref{cor:quotient-minimal} says that this is the smallest possible
nonzero layer on which the obstruction can live.  The coarser quotient by
\(I=\m^3\) makes the doubling law transparent; the finer quotient by the
socle isolates precisely the obstruction class.

For comparison, there are square-zero quotients internal to the Hopf
algebra.  The ideal \((\theta)=(\lambda-1)\) is a square-zero Hopf ideal,
because
\[
 \Delta(\theta)=\theta\ot1+1\ot\theta+\theta\ot\theta,
 \qquad S(\theta)=-\theta.
\]
Modulo it, \(\lambda=1\), so \(U,V\) are primitive and the quotient is
killed by four.  Likewise \((cV)\) is a square-zero Hopf ideal by
\eqref{eq:cV-primitive}.  These quotients are not flat over \(C\), so they
do not describe deformations through rank-four finite flat group schemes.
It is also important that the final defect ideal \((cW)=\Soc(A)\) is not a
Hopf ideal: direct expansion gives
\[
 \Delta(cW)=c\bigl(W\ot1+U\ot V+V\ot U+1\ot W\bigr).
\]
Thus the primitive source \(cV\), rather than only its product \(cW\), is
the intrinsic Hopf-theoretic obstruction to kill.

There is also an elementary cyclic-cohomological reading of the norm.  For
a \(C\)-algebra \(R\) and \(g\in G(R)\), put
\(a=\lambda(g)\) and \(u=U(g)\).  The skew-primitive product law says that
the restriction of \(U\) along \(\Z\to G(R)\), \(n\mapsto g^n\), is the
crossed homomorphism, for \(n\ge0\),
\[
 z_g(n)=N_n(a)u.
\]
Thus \(z_g\) is a \(1\)-cocycle for the action of \(\Z\) on \(R\) in which
one acts by \(a\).  Since \(a^4=1\), the action factors through the cyclic
group \(C_4\).  The cocycle descends to \(C_4\) precisely when its cyclic
norm
\[
 z_g(4)=N_4(a)u
\]
vanishes.  Universally this norm is \(2sW\), so the cocycle does not descend
to \(C_4\); it does descend to \(C_8\), because \(N_8(\lambda)=0\).  This
is exactly the norm differential in the standard periodic resolution of a
cyclic group.  Notice the scope of the statement: it concerns the
one-generated subgroup attached to each \(g\).  Since \(G\) is
noncommutative, the global power map \([4]\) is not itself being claimed to
be a group homomorphism or an ordinary group-cohomology class.

The two square-zero quotients fit into
\[
 C\longrightarrow\overline C=C/(2s)
 \longrightarrow B=C/\m^3.
\]
The ring \(B\) has characteristic two.  Lifting from \(B\) to
\(\overline C\) restores characteristic four but still does not break the
fourth-power relation.  Only the final one-dimensional socle extension
from \(\overline C\) to \(C\) creates the obstruction.  Altogether the
nilpotent deformation raises the exponent in two visible steps:
\[
 \begin{array}{c|c|c|c}
  \text{base}&\text{characteristic}&\text{power behavior}
             &\text{exact bound}\\ \hline
  \F_2&2&[2]=\unit&2\\
  B=\F_2[r,s]/(r,s)^3&2&[2]\ne\unit,\ [4]=\unit&4\\
  \overline C=C/(2s)&4&[2]\ne\unit,\ [4]=\unit&4\\
  C&4&[4]\ne\unit,\ [8]=\unit&8.
 \end{array}
\]
Thus it is the group-theoretic equation \([4]=\unit\), rather than the
scalar equation \(4=0\), which breaks under the square-zero deformation:
the final base ring still has characteristic four.

\section{The closed fiber and the total family}\label{sec:geometry}

The formulas become more informative when compared with the closed fiber.
Modulo \(\m=(r,s)\),
\[
 A\ot_C\F_2=\F_2[U,V]/(U^2,V^2),
 \qquad \lambda=1,
\]
and both \(U\) and \(V\) are primitive.  Hence
\begin{equation}\label{eq:special-fiber}
 G_{\F_2}\cong\alpha _2\times\alpha _2.
\end{equation}
The special fiber is commutative and is already killed by two.  Thus no
geometric point witnesses the failure of killedness; the fourth-power word
is a genuinely vertical morphism supported in the nilpotent structure of
the base.

The total group scheme is not commutative.  Let
\(\tau\colon A\ot_C A\to A\ot_C A\) exchange the two factors.  From
\eqref{eq:Delta-main},
\begin{equation}\label{eq:noncocommutative-defect}
 \Delta(U)-\tau\Delta(U)
 =\theta\ot U-U\ot\theta.
\end{equation}
The coefficient of \(V\ot U\) is \(s\), which is nonzero because
\(2s\ne0\).  Since \(A\ot_C A\) is free on the sixteen tensor products of
\(1,U,V,W\), the right side of
\eqref{eq:noncocommutative-defect} is nonzero.  Thus the coproduct is not
cocommutative and \(G\) is noncommutative.  Both commutativity and the
power relation therefore jump between the closed fiber and the total
family, as the positive results recalled in the introduction predict.

There is also a useful rank-two shadow of the construction.  The equation
\(V=0\) defines a finite locally free closed subgroup
\begin{equation}\label{eq:H0}
 H_0\subset G,
 \qquad
 \mathcal O(H_0)=C[U]/(U^2-r^2U).
\end{equation}
Indeed, \((V)\) is a Hopf ideal.  On \(H_0\), the character is
\(1+rU\), and the rank-two calculation gives
\[
 [2]^\#(U)=(2+rU)U=(2+r^3)U=0.
\]
Thus \(H_0\) is killed by two.  Since \([4]^\#(V)=0\), the fourth-power
word of \(G\) factors nontrivially through \(H_0\); squaring once more
kills its image.  This gives a group-theoretic complement to the norm
explanation of \([8]=\unit\).

The subgroup \(H_0\) is not normal.  To see the obstruction directly, let
\(g=(u,v,z)\) be a point of \(G\), and let
\(h=(a,0,1+ra)\) be a point of \(H_0\).  Conjugation inside the ambient
group \eqref{eq:H-law} gives
\begin{equation}\label{eq:conjugation}
 v(ghg^{-1})=rva z^{-1}.
\end{equation}
Universally, its coefficient is
\[
 rV\lambda^{-1}=rV-r^2W\ne0.
\]
Indeed \(r\ne0\), since \(r^3=-2\) and \(2s\ne0\), and freeness of \(A\)
proves the displayed nonvanishing.  The coordinate \(a\) is part of a free
\(C\)-basis of \(\mathcal O(H_0)\), so the universal conjugation coordinate
is nonzero as well.
Hence the conjugation morphism does not factor through \(V=0\).  The two
rank-two patterns should therefore not be interpreted as a normal
extension of one rank-two group by another: the bridge destroys precisely
that naive decomposition.

\section{Economy and scope}\label{sec:economy}

The presentation has two distinct and useful minimality properties.  The
first was proved in Corollary~\ref{cor:quotient-minimal}: every proper
quotient by a nonzero ideal kills \(2s\), hence kills the displayed
fourth-power obstruction.  Thus the witness is supported on the final
one-dimensional socle extension of its base.

The second concerns the number of equations in the chosen parameter
space.  Let
\[
 I=(r^3+2,s^3+2,rs^2)\subset\Z[r,s],
 \qquad
 \mathfrak p=(2,r,s).
\]
The consequences \(4=r^4=s^5=0\) in \(C\) show that
\[
 \sqrt I=\mathfrak p.
\]
The maximal ideal \(\mathfrak p\) has height three.  Krull's height theorem
therefore implies that \(I\) cannot be generated by fewer than three
elements.  More generally, every proper ideal \(J\supset I\) has radical
\(\mathfrak p\) and also requires at least three generators.  Thus the
closed compatibility locus used in this paper cannot be cut out inside
\(\Spec\Z[r,s]\) by fewer equations.

This statement has a deliberately restricted scope.  The number of
relations in an arbitrary presentation is not an intrinsic invariant, and
we do not claim that every possible rank-four counterexample must use
three equations or a length-nine base.  What the argument proves is
precise: after choosing the two parameters \(r,s\) and the bridge
coefficients
\[
 \alpha=r^2,
 \qquad \beta=-rs,
 \qquad \gamma=s^2,
\]
the two diagonal cancellations and the one mixed vanishing give an
equation-minimal locus in this fixed ambient space.  No fourth relation and
no independent Artinian truncation is hidden in the construction.

\section{Perspective}

The example can be remembered through the following chain:
\[
 \boxed{
 \begin{gathered}
 \text{two rank-two cancellations and one bridge}\\
 \Downarrow\\
 \text{a nearly trivial character }\lambda=1+\theta,
 \quad\theta^2=0\\
 \Downarrow\\
 N_4(\lambda)=2\theta=2sV\\
 \Downarrow\\
 [4]^\#(U)=2sUV\in\Soc(A),\qquad [8]=\unit.
 \end{gathered}}
\]
The main point is not the particular reduction table.  It is that a
geometric sum in characteristic four can remember a square-zero
deformation which all geometric fibers forget.  Skew-primitive coordinates
turn that arithmetic carry into a power word, and the bridge moves it to
the last infinitesimal layer where it cannot cancel prematurely.

This viewpoint suggests a broader way to look for relative power
obstructions: begin with a group-like element infinitesimally close to one,
study the relevant geometric sum, and then ask for a finite flat algebra
which stores the surviving coefficient in its socle.  In the present
example all of those roles are visible in two coordinates and three base
relations.

\appendix

\section{The subgroup calculation}\label{app:closure}

We record the reductions used in Proposition~\ref{prop:submonoid}.  This
appendix contains the bookkeeping so that the main argument can emphasize
the ambient group.

Four consequences of the defining relations are
\begin{align}
 s^2(\lambda-1)&=-2V,\label{eq:aux-1}\\
 rs(\lambda-1)&=r^2sU,\label{eq:aux-2}\\
 rs(\lambda^2-1)&=0,\label{eq:aux-3}\\
 2\lambda U+r^2(\lambda^2-\lambda)+r^2sV&=0.
 \label{eq:aux-4}
\end{align}
For example, the first two follow by multiplying
\(\lambda-1=rU+sV+rsW\) by \(s^2\) and \(rs\), respectively.  The third
uses \(\lambda^2-1=2sV\).  For the fourth, note that
\(\lambda^2-\lambda=\theta\), and expand using \(2\theta=2sV\).

For \(V_*=V_1\lambda_2+V_2\), one obtains
\begin{align*}
 V_*^2-s^2V_*
 &=V_1\lambda_2\bigl(s^2(\lambda_2-1)+2V_2\bigr)=0
\end{align*}
by \eqref{eq:aux-1}.  For \(U_*=U_1+\lambda_1U_2\), expansion gives
\begin{align*}
 U_*^2-r^2U_*+rsV_*
 &=-rsV_1(1-\lambda_2)\\
 &\quad+
 \bigl(2\lambda_1U_1+r^2(\lambda_1^2-\lambda_1)\bigr)U_2
 -rs(\lambda_1^2-1)V_2.
\end{align*}
Equations \eqref{eq:aux-2} and \eqref{eq:aux-3} turn this into
\[
 \bigl(r^2sV_1+2\lambda_1U_1
       +r^2(\lambda_1^2-\lambda_1)\bigr)U_2,
\]
which vanishes by \eqref{eq:aux-4}.  This proves preservation of both
quadrics.  The graph equation is the factorization
\eqref{eq:graph-factor} proved in the main text.

As an optional independent check, the ambient inverse is
\[
 (u,v,z)^{-1}=(-z^{-1}u,-vz^{-1},z^{-1}).
\]
Substitution of \eqref{eq:explicit-antipode} into the two quadrics gives
zero, and \(S(\lambda)=\lambda^{-1}\).  Thus one may also prove inverse
stability directly, although Proposition~\ref{prop:inverse-by-seven} makes
that calculation unnecessary.

\section{Exact algebra audit}\label{app:audit}

All proofs above are by hand.  For reproducibility, the polynomial
identities were also checked by exact strong Gr\"obner reduction over
\(\Z\).  With lexicographic order \(r>s\), a strong Gr\"obner basis of the
base ideal is
\[
 (4,\ 2s^2,\ s^3+2,\ 2r,\ rs^2,\ r^3+2).
\]
It leaves \(2s\) in nonzero normal form.  The accompanying bounded
Macaulay2 scripts are
\begin{center}
 \path{m2/verify_rank4_counterexample_simple.m2},
 \qquad
 \path{m2/audit_rank4_unoptimized_bridge_20260711.m2}.
\end{center}
From the project root it can be run with
\begin{verbatim}
M2 --script m2/verify_rank4_counterexample_simple.m2
M2 --script m2/audit_rank4_unoptimized_bridge_20260711.m2
\end{verbatim}
The first script independently verifies the regular-sequence ideal quotients,
the base relations and Loewy filtration, both socles, the tangent-cone
Hilbert function, and the special fiber.  It also checks the elementary
identities of Lemma~\ref{lem:theta}, the four reductions in
Appendix~\ref{app:closure}, and the graph factorization
\eqref{eq:graph-factor}.

In particular, the Hopf algebra assertion is checked at the level of the
presentation, not merely inferred from the power calculation.  The script
verifies that the proposed coproduct preserves both defining quadrics and
that \(\Delta(\lambda)=\lambda\ot\lambda\); it checks both counit
identities on \(U,V\); it compares the two parenthesizations of the product
on three universal points; and it verifies that the displayed antipode
preserves the relations and satisfies both convolution identities on
\(U,V\).  Since \(A\) is generated as a \(C\)-algebra by \(U,V\), these
are precisely the bialgebra and antipode axioms.  The same computation
checks noncocommutativity, the rank-two Hopf ideal and its nonnormality,
agreement of the seventh-power word with the antipode, and the four normal
forms
\[
 \bigl([4]^\#(U),[4]^\#(V),[8]^\#(U),[8]^\#(V)\bigr)
 =(2sUV,0,0,0).
\]
It further verifies the square-zero identities for \((rs)\), \(\m^3\),
and \((2s)\); the two character cocycles
\eqref{eq:character-cocycles}; the first-order coproduct
\eqref{eq:first-order-coproduct}; the reduced doubling formulas
\eqref{eq:doubling-downstairs}; and the vanishing of the fourth-power word
after each relevant base quotient.  It also checks the two internal Hopf
quotients above and verifies directly that the socle defect ideal itself is
not a Hopf ideal.  Finally, it verifies the naive-lift curvature
\eqref{eq:naive-lift-curvature}.  The second script works over the
untruncated five-parameter bridge ring: it verifies that
\eqref{eq:bridge-chart-again} preserves both quadrics, makes \(\lambda\)
group-like, supplies the displayed antipode, and gives
\([4]^\#(U)=2sUV\).  It then checks all eight framed lifts in
\eqref{eq:eight-triangular-lifts} separately and confirms that every one
has the same fourth-power carry.
The remaining passages in the paper are formal deductions rather than
additional polynomial identities: for example, monic division gives the
free bases and the lengths, the regular sequences give the complete
intersection assertions, and the two minimality statements use the
standard socle and height arguments written in the text.  Historical and
literature statements are of course outside the scope of a computer
algebra audit.  Thus the computation checks every algebraic identity used
in the construction and the explicit inputs to the structural arguments;
it remains supplementary, and no proof in the paper depends on it.

\begin{thebibliography}{99}

\bibitem{SGA3I}
M.~Demazure and A.~Grothendieck (eds.),
\emph{Sch\'emas en groupes I: Propri\'et\'es g\'en\'erales des sch\'emas
en groupes (SGA 3)},
Lecture Notes in Mathematics, vol.~151, Springer-Verlag, 1970.

\bibitem{SGA3II}
M.~Demazure and A.~Grothendieck (eds.),
\emph{Sch\'emas en groupes II: Groupes de type multiplicatif, et structure
des sch\'emas en groupes g\'en\'eraux (SGA 3)},
Lecture Notes in Mathematics, vol.~152, Springer-Verlag, 1970.

\bibitem{GerstenhaberSchack}
M.~Gerstenhaber and S.~D.~Schack,
\emph{Bialgebra cohomology, deformations, and quantum groups},
Proc. Natl. Acad. Sci. USA \textbf{87} (1990), no.~1, 478--481,
\href{https://doi.org/10.1073/pnas.87.1.478}
{doi:10.1073/pnas.87.1.478}.

\bibitem{IllusieCotangent}
L.~Illusie,
\emph{Complexe cotangent et d\'eformations I},
Lecture Notes in Mathematics, vol.~239, Springer-Verlag, 1971,
\href{https://doi.org/10.1007/BFb0059052}
{doi:10.1007/BFb0059052}.

\bibitem{OortTate}
J.~Tate and F.~Oort,
\emph{Group schemes of prime order},
Ann. Sci. \'Ecole Norm. Sup. (4) \textbf{3} (1970), no.~1, 1--21.

\bibitem{Schoof2001}
R.~Schoof,
\emph{Finite flat group schemes over local Artin rings},
Compos. Math. \textbf{128} (2001), no.~1, 1--15,
\href{https://doi.org/10.1023/A:1017560215203}{doi:10.1023/A:1017560215203}.

\bibitem{SchoofSurvey}
R.~Schoof,
\emph{Is a finite locally free group scheme killed by its order?},
in \emph{Open Problems in Arithmetic Algebraic Geometry},
Advanced Lectures in Mathematics, vol.~46, International Press, 2019,
1--7.

\bibitem{Stix}
J.~Stix,
\emph{A course on finite flat group schemes and \(p\)-divisible groups},
lecture notes, Heidelberg, 2009, revised 2012,
\href{https://www.math.uni-frankfurt.de/~stix/skripte/STIXfinflatGrpschemes20120918.pdf}
{available online}.

\bibitem{TateFiniteFlat}
J.~Tate,
\emph{Finite flat group schemes},
in \emph{Modular Forms and Fermat's Last Theorem},
G.~Cornell, J.~H.~Silverman, and G.~Stevens (eds.),
Springer-Verlag, New York, 1997, 121--154.

\bibitem{Torti}
E.~Torti,
\emph{Lagrange's theorem for a class of finite flat group schemes over
local Artin rings},
Doc. Math. (2025), online first,
\href{https://doi.org/10.4171/DM/1055}{doi:10.4171/DM/1055}.

\end{thebibliography}

\end{document}
